\bibitem[Ilkoğlu and Macedo(2024)]{ilkoglu2024}
Ilkoğlu, E., Macedo, J. (2024). A social equity analysis of public EV charging access in the US.
\textit{Energy Research \& Social Science}.
\doi{10.1016/j.erss.2024.103301}

\bibitem[Hecht et~al.(2022)]{hecht2022}
Hecht, C., Figgener, J., Sauer, D.~U. (2022). Empirical analysis of public charging infrastructure utilization and profitability in Germany.
\textit{arXiv preprint arXiv:2206.09582}.

\bibitem[Liu et~al.(2021)]{liu2021}
Liu, S., Lin, C., Ouyang, Y. (2021). A network equilibrium framework for electric vehicle charging station location planning.
\textit{arXiv preprint arXiv:2102.05851}.

\bibitem[Magsino et~al.(2024)]{magsino2024}
Magsino, L., Adkins, E., Mooney, P. (2024). Optimizing EV charging infrastructure using spatial clustering of mobility data.
\textit{ISPRS International Journal of Geo-Information}.
\doi{10.3390/ijgi13100368}

\bibitem[Pan et~al.(2024)]{pan2024}
Pan, S., Liu, X., Wang, Y. (2024). Built environment, public charging access, and electric vehicle adoption in the US: A spatial regression analysis.
\textit{arXiv preprint arXiv:2401.17456}.

\bibitem[Çelik and Ok(2024)]{celik2024}
Çelik, H., Ok, E. (2024). P-median model for EV fast charger deployment under demand uncertainty.
\textit{Heliyon}.
\doi{10.1016/j.heliyon.2024.e11055}

\bibitem[Padmanabhan et~al.(2021)]{padmanabhan2021}
Padmanabhan, R., Wang, X., Wang, F. (2021). Location optimization of EV charging stations using reinforcement learning.
\textit{Transportation Research Part A: Policy and Practice}.
\doi{10.1016/j.tra.2021.08.001}

\bibitem[Kuang et~al.(2024)]{kuang2024}
Kuang et~al. (2024). CityEVCP: Context-aware EV charging prediction and planning using hypergraph attention networks.
\textit{arXiv preprint arXiv:2401.18766}.

\bibitem[Wang et~al.(2021)]{wang2021}
Wang et~al. (2021). SPAP: Simultaneous prediction and planning framework for electric vehicle charging infrastructure.
\textit{arXiv preprint arXiv:2110.09452}.

\bibitem[Tra et~al.(2022)]{tra2022}
Tra, N., Carvalho, A., Silva, M. (2022). Max-covering and user-choice hybrid models for electric vehicle charging station location.
\textit{arXiv preprint arXiv:2206.11165}.

\bibitem[Zhang and Liu(2022)]{hongkong2022}
Zhang, Y., Liu, W. (2022). Optimal allocation of electric vehicle charging stations in Hong Kong using discrete choice modeling.
\textit{Transportation Research Part A: Policy and Practice}.
\doi{10.1016/j.tra.2022.02.005}

\bibitem[Dharangaonkar and Patil(2025)]{dharangaonkar2025}
Dharangaonkar, P., Patil, A. (2025). Systematic review of electric vehicle charging infrastructure planning models: Gaps and future directions.
\textit{International Journal of Sustainable Development and Planning}.
\doi{10.18280/ijsdp.200605}

\bibitem[Agency(2025)]{iea2025}
International Energy Agency (2025). Global EV Outlook 2025.
\url{https://www.iea.org/reports/global-ev-outlook-2025}. Accessed July 2025.

\bibitem[Consultancy(2024)]{guardian2024}
Consultancy, S. (2024). UK faces severe EV charger access inequality, study shows.
\textit{The Guardian}.
Accessed: July 2025.

\bibitem[Zhou et~al.(2018)]{etrr2018}
Zhou, Y., Levin, J., Zhang, L. (2018). Regional disparities in EV public charging infrastructure: Evidence from California.
\textit{European Transport Research Review}.
\doi{10.1186/s12544-018-0322-8}

\bibitem[Carlton and Sultana(2022)]{carlton2022}
Carlton, G.~J., Sultana, S. (2022). Electric vehicle charging station accessibility and land use clustering: A case study of the Chicago region.
\textit{Journal of Urban Mobility}.
\doi{https://doi.org/10.1016/j.urbmob.2022.100019}

\bibitem[Qu et~al.(2024)]{qu2024}
Qu et~al. (2024). A Physics-Informed and Attention-Based Graph Learning Approach for Regional Electric Vehicle Charging Demand Prediction.
\textit{IEEE Transactions on Intelligent Transportation Systems}.
\doi{10.1109/TITS.2024.3401850}

\bibitem[Torkey and Abdelgawad(2022)]{torkey2022}
Torkey, A., Abdelgawad, H. (2022). Framework for planning of EV charging infrastructure: Where should cities start?.
\textit{Transport Policy}.

\bibitem[Wong(2004)]{wong2004modifiable}
Wong, D.~W. (2004). The modifiable areal unit problem (MAUP).
\textit{WorldMinds: geographical perspectives on 100 problems: commemorating the 100th anniversary of the association of American geographers 1904--2004}.

\bibitem[IEA(2023)]{IEA2023GlobalEVOutlook}
International Energy Agency (IEA) (2023). Global EV Outlook 2023: Catching Up with Climate Ambitions.
\url{https://www.iea.org/reports/global-ev-outlook-2023}. License: CC BY 4.0.

\bibitem[Ghosh et~al.(2022)]{Ghosh2022}
Ghosh, N., Mothilal Bhagavathy, S., Thakur, J. (2022). Accelerating electric vehicle adoption: techno-economic assessment to modify existing fuel stations with fast charging infrastructure.
\textit{Clean Technologies and Environmental Policy}.
\doi{10.1007/s10098-022-02406-x}

\bibitem[Cai et~al.(2024)]{Cai2024}
Cai et~al. (2024). An Analytical Framework for Assessing Equity of Access to Public Electric Vehicle Charging Stations: The Case of Shanghai.
\textit{Sustainability}.
\doi{10.3390/su16146196}

\bibitem[Ullah et~al.(2023)]{Ullah2023}
Ullah et~al. (2023). Optimal Deployment of Electric Vehicles’ Fast-Charging Stations.
\textit{Journal of Advanced Transportation}.
\doi{10.1155/2023/6103796}

\bibitem[Ullah et~al.(2023)]{ullah2023}
Ullah et~al. (2023). Optimal Deployment of Electric Vehicles’ Fast-Charging Stations.
\textit{Journal of Advanced Transportation}.
\doi{10.1155/2023/6103796}

\bibitem[Heo and Chang(2024)]{HEO2024}
Heo, J., Chang, S. (2024). Optimal planning for electric vehicle fast charging stations placements in a city scale using an advantage actor-critic deep reinforcement learning and geospatial analysis.
\textit{Sustainable Cities and Society}.
\doi{https://doi.org/10.1016/j.scs.2024.105567}

\bibitem[Kavianipour et~al.(2021)]{KAVIANIPOUR2021}
Kavianipour et~al. (2021). Electric vehicle fast charging infrastructure planning in urban networks considering daily travel and charging behavior.
\textit{Transportation Research Part D: Transport and Environment}.
\doi{https://doi.org/10.1016/j.trd.2021.102769}

\bibitem[E. Seilabi et~al.(2025)]{Seilabi2025}
E. Seilabi et~al. (2025). Sustainable Planning of Electric Vehicle Charging Stations: A Bi-Level Optimization Framework for Reducing Vehicular Emissions in Urban Road Networks.
\textit{Sustainability}.
\doi{10.3390/su17010001}

\bibitem[An et~al.(2025)]{AN2025}
An et~al. (2025). Urban public charging infrastructure planning for electric vehicles: A continuum approximation approach.
\textit{Transportation Research Part D: Transport and Environment}.
\doi{https://doi.org/10.1016/j.trd.2025.104605}

\bibitem[Karmaker et~al.(2025)]{KARMAKER2025}
Karmaker et~al. (2025). Customer-centric meso-level planning for electric vehicle charger distribution.
\textit{Applied Energy}.
\doi{https://doi.org/10.1016/j.apenergy.2025.125742}

\bibitem[Kapoor et~al.(2024)]{Kapoor2024}
Kapoor et~al. (2024). Optimal Planning of Fast EV Charging Stations in a Coupled Transportation and Electrical Power Distribution Network.
\textit{IEEE Transactions on Automation Science and Engineering}.
\doi{10.1109/TASE.2023.3293955}

\bibitem[Pal et~al.(2023)]{10018877}
Pal, A., Bhattacharya, A., Chakraborty, A.~K. (2023). Planning of EV Charging Station With Distribution Network Expansion Considering Traffic Congestion and Uncertainties.
\textit{IEEE Transactions on Industry Applications}.
\doi{10.1109/TIA.2023.3237650}

\bibitem[Jha et~al.(2025)]{JHA2025}
Jha et~al. (2025). An equity-based approach for addressing inequality in electric vehicle charging infrastructure: Leaving no one behind in transport electrification.
\textit{Energy for Sustainable Development}.
\doi{https://doi.org/10.1016/j.esd.2024.101643}

\bibitem[Iravani(2022)]{IRAVANI2022}
Iravani, H. (2022). A multicriteria GIS-based decision-making approach for locating electric vehicle charging stations.
\textit{Transportation Engineering}.
\doi{https://doi.org/10.1016/j.treng.2022.100135}

\bibitem[Roy and Law(2022)]{ROY2022}
Roy, A., Law, M. (2022). Examining spatial disparities in electric vehicle charging station placements using machine learning.
\textit{Sustainable Cities and Society}.
\doi{https://doi.org/10.1016/j.scs.2022.103978}

\bibitem[Law and Roy(2021)]{Law2021}
Law, M., Roy, A. (2021). A geospatial data fusion framework to quantify variations in electric vehicle charging demand.
\textit{Proceedings of the 4th ACM SIGSPATIAL International Workshop on Advances in Resilient and Intelligent Cities}.
\doi{10.1145/3486626.3493429}

\bibitem[Heo and Chang(2024)]{heo_optimal_2024}
Heo, J., Chang, S. (2024). Optimal planning for electric vehicle fast charging stations placements in a city scale using an advantage actor-critic deep reinforcement learning and geospatial analysis.
\textit{Sustainable Cities and Society}.
\doi{10.1016/j.scs.2024.105567}

\bibitem[Bian et~al.(2022)]{bian_planning_2022}
Bian et~al. (2022). Planning of electric vehicle fast-charging station based on POI.
\textit{Energy Reports}.
\doi{https://doi.org/10.1016/j.egyr.2022.10.161}

\bibitem[Shi(2018)]{shi2018}
Shi, C. (2018). Propulsion-Machine-Integrated Universal Onboard Chargers for Electric Vehicles.
\doi{10.13016/M2CN6Z32S}

\bibitem[Luo et~al.(2024)]{Luo2024}
Luo et~al. (2024). Innovation diffusion in EV charging infrastructure.
\textit{Transportation Research Part C: Emerging Technologies}.
\doi{https://doi.org/10.1016/j.trc.2024.104733}

\bibitem[Metais et~al.(2022)]{Metais2022}
Metais et~al. (2022). Too much or not enough? Planning EV charging infrastructure.
\textit{Renewable and Sustainable Energy Reviews}.
\doi{10.1016/j.rser.2021.111719}

\bibitem[Statistics Canada(2011a)]{StatCan_CT_2011_Def}
Statistics Canada (2011). Census tract (CT).
\url{https://www150.statcan.gc.ca/n1/pub/92-195-x/2011001/geo/ct-sr/def-eng.htm}.
From the 2011 Census Dictionary, Catalogue no.\ 92-195-X.

\bibitem[Statistics Canada(2021)]{StatCan_FSA_2021_Def}
Statistics Canada (2021). Forward sortation area (FSA).
\url{https://www12.statcan.gc.ca/census-recensement/2021/ref/dict/az/Definition-eng.cfm?ID=geo036}.
From the 2021 Census Dictionary, Catalogue no.\ 98-301-X.

\bibitem[Statistics Canada(2011b)]{StatCan_DA_2011_Def}
Statistics Canada (2011). Dissemination area (DA).
\url{https://www150.statcan.gc.ca/n1/pub/92-195-x/2011001/geo/da-ad/def-eng.htm}.
From the 2011 Census Dictionary, Catalogue no.\ 92-195-X.

